\documentclass{article}
\usepackage[utf8]{inputenc}
\usepackage{amsmath}
\usepackage{geometry}
\geometry{margin=2cm}
\begin{document}

\section*{Questão 124 — ENEM 2024 — Ciências da Natureza}

\subsection*{Tema: Acústica e Inteligibilidade Sonora}

Esta questão avalia a compreensão sobre o impacto do nível sonoro da fala e do ruído de fundo na inteligibilidade em ambientes fechados, especialmente em uma sala de aula.

\subsection*{Interpretação e Estratégia}
Considerando que o professor emite um som de $60\text{ dB}$ a 1 m e que o ruído de fundo pode chegar a $45\text{ dB}$, sabe-se que a fala deve estar pelo menos $5\text{ dB}$ acima do ruído para ser compreendida. Assim, o sinal mínimo aceitável é $50\text{ dB}$.

O gráfico fornecido relaciona o nível sonoro à distância, permitindo identificar a maior distância onde a fala permanece inteligível.

\subsection*{Análise do Conteúdo e Mini Revisão}
\begin{table}[h!]
\centering
\begin{tabular}{|p{5cm}|p{6cm}|p{6cm}|}
\hline
\textbf{Conceito} & \textbf{Ideia-chave} & \textbf{Aplicação} \\
\hline
Decaimento do nível sonoro com a distância & O som perde intensidade conforme a distância da fonte aumenta & Determinar alcance máximo para compreensão da fala \\
\hline
Relação entre fala e ruído de fundo & A fala deve exceder o ruído de fundo em pelo menos 5 dB & Define o nível sonoro mínimo para a inteligibilidade \\
\hline
Interpretação do gráfico & Leitura correta do nível sonoro versus distância & Localizar a distância máxima para compreensão da fala \\
\hline
\end{tabular}
\caption{Revisão rápida dos conceitos usados}
\end{table}

\subsection*{Resolução e Pulo do Gato}
\begin{itemize}
  \item Nível sonoro da fala a 1 m: $60\text{ dB}$.
  \item Ruído máximo: $45\text{ dB}$.
  \item Nível mínimo para entender a fala: $45 + 5 = 50\text{ dB}$.
  \item Pelo gráfico, o nível sonoro cai para cerca de $50\text{ dB}$ aos 3 metros.
  \item Portanto, a maior distância para compreensão é aproximadamente $3,0$ metros.
\end{itemize}

\paragraph{Pulo do gato:} Para garantir que a fala seja compreendida, o nível sonoro deve estar pelo menos $5\text{ dB}$ acima do ruído de fundo. Use o gráfico para achar essa distância exata, observando onde a curva atinge o nível mínimo.

\subsection*{Armadilhas na resolução}
\begin{itemize}
  \item Confundir o valor do nível sonoro necessário para compreensão.
  \item Interpretação incorreta do gráfico levando a superestimar distâncias.
  \item Ignorar o ruído de fundo na análise.
  \item Suposições equivocadas sobre a diminuição do som.
\end{itemize}

\subsection*{Código da questão}
CN\_ACUSTICA\_INTERPRETACAO\_001

\bigskip

\textbf{Resposta:} $\boxed{3,0\text{ m}}$ (alternativa correta).

\end{document}